\chapter{Desarrollo e Implementación de la Solución}\label{desarrollo}

\azul{Este capítulo documenta la fase ejecutiva e ingenieril del proyecto, detallando la construcción, modelado, cálculo o programación del artefacto propuesto. En este apartado se explicita la transformación de los requerimientos y diseños previos en un producto, sistema o estructura tangible. La literatura en ingeniería sostiene que la documentación detallada del desarrollo garantiza la reproducibilidad y mantenibilidad de la solución \cite[p.~215]{pressman2010ingenieria}.}

\section{Selección de Herramientas, Materiales o Lenguajes}

\azul{Esta sección fundamenta técnicamente las herramientas, insumos o tecnologías seleccionadas para materializar la propuesta. Su alcance varía según la especialidad de la ingeniería:}

\begin{itemize}
	\item \azul{\textbf{Exclusivo de Ingeniería Civil en Informática:} Se justifica la elección del lenguaje de programación principal, marcos de trabajo (\textit{frameworks}), motores de bases de datos y bibliotecas, analizando sus propiedades computacionales, paradigma y gestión de memoria.}
	\item \azul{\textbf{Otras Ingenierías Civiles (Obras Civiles, Industrial, Mecánica, etc.):} Esta sección se adapta para justificar la selección de materiales de construcción, componentes mecánicos/eléctricos, instrumental de medición, software de simulación estandarizado (ej. SAP2000, ANSYS, FlexSim) o normativas técnicas aplicadas.}
\end{itemize}

\subsection{Propiedades Técnicas y Justificación}

\azul{Se analizan las características intrínsecas de las herramientas o materiales seleccionados que justifican su adopción frente al problema abordado.}

\moradoc{Ejemplo (Informática): Se seleccionó Python debido a su soporte nativo para bibliotecas de manipulación de datos y su integración eficiente con motores relacionales mediante ORM.}

\moradoc{Ejemplo (Obras Civiles / Mecánica): Se seleccionó acero de alta resistencia grado A630-420H debido a sus propiedades de ductilidad y comportamiento ante solicitaciones sísmicas en la estructura.}

\subsection{Ventajas y Desventajas Comparativas}

\azul{Se presenta un análisis crítico que evalúa los beneficios e inconvenientes de las opciones seleccionadas en comparación con otras alternativas del mercado o la industria.}

\section{Arquitectura y Diseño Detallado de la Solución}

\azul{Se describe la estructura global del sistema o artefacto y la interacción entre sus componentes. En \textbf{Ingeniería Civil en Informática}, este apartado se fundamenta en diagramas de arquitectura de software (UML, modelos de capas o microservicios). En \textbf{otras especialidades}, comprende diagramas de flujo de procesos (P\&ID), esquemas eléctricos, planos estructurales o diagramas de planta.}

\section{Construcción e Integración de Componentes}

\azul{Describe el procedimiento paso a paso mediante el cual se ensamblaron, codificaron o erigieron las distintas partes del proyecto. Se sugiere acompañar el texto con fragmentos de algoritmos representativos, esquemas de montaje o memorias de cálculo.}

